本文围绕\textbf{带相变材料的低温防护服}，采用拓展分离变量法等方法，研究不同环境、结构和材料性能下防护服的有效防护时间。

\textbf{针对问题一，本文建立基于有效热容法的含相变层三层非稳态导热模型。} 首先对 DSC（差示扫描量热）曲线进行基线校正，并由累计放热量构造功能层不可逆固化进度，将其引入有效热容，建立三层一维非稳态导热方程。采用\textbf{拓展分离变量法}分时段求解温度场。计算得到，进入 $-40\,^{\circ}\mathrm C$ 环境 \textbf{273.91 秒}后人体表面温度达到 \textbf{$15\,^{\circ}\mathrm C$}，工作将很困难；\textbf{311.16 秒}后温度降至 \textbf{$10\,^{\circ}\mathrm C$}，将有生命危险。

\textbf{针对问题二，本文进一步建立有风轻微运动条件下的三层非稳态导热模型。} 首先根据环境风速和轻微运动状态确定防护服内外换热参数，分别建立\textbf{人体与防护服内表面}和\textbf{防护服外表面与有风环境的对流 Robin 边界}。在此基础上形成有风轻微运动条件下的多层非稳态热传导初边值问题。计算得到，\textbf{103.49 秒}后，工作将很困难；\textbf{115.67 秒}后，将有生命危险。相比静止无风环境，工作困难时刻提前约 \textbf{170.42 秒}。

\textbf{针对问题三，本文进一步建立以最大有效持续时间为目标的 PDE 约束离散优化模型。} 首先根据材料费用上限确定隔热层可行增加厚度，并将随站立时间衰减的承重能力转化为约束条件。将不同厚度方案的热防护时间与人体承重限制共同纳入优化。计算得到最优方案为新增三层外涂层，最外隔热层厚度为 \textbf{1.2 毫米}、防护服总厚度为 \textbf{2.3 毫米}；\textbf{416.04 秒}后人体工作将很困难，\textbf{490.64 秒}后将有生命危险。与原防护服相比，两种人体状态对应的有效防护时间分别提高约 \textbf{51.89\% 和 57.68\%}。

\textbf{针对问题四，本文进一步建立以最小放热强化率为目标的相变材料热传导参数反问题模型。} 首先将功能层各温度下的放热能力按统一比例增强，并保持相变温区及固化进度分布不变。以求得的\textbf{最大有效持续时间}为目标热响应，反求功能层单位质量相变潜热所需的最小强化程度。采用\textbf{上界扩张与二分搜索法}求解。计算得到，为达到\textbf{人体工作困难前和生命危险前的最大有效持续时间}，功能层放热能力至少需要提高 \textbf{86.33\% 和 108.98\%}。

\textbf{本文特色在于：} 基于 DSC 曲线累计放热量构造不可逆固化进度。

\keywords{低温防护服；相变材料；有效热容法；拓展分离变量法；偏微分方程约束优化；热传导参数反问题}
